\documentclass{article}
\usepackage[utf8]{inputenc}
\usepackage[russian]{babel} % For Russian language support
\usepackage{amsmath}
\usepackage{amsfonts}
\usepackage{amssymb}

\title{Пример научной статьи}
\author{Автор}
\date{\today}

\begin{document}

\maketitle

\section{Введение}
Настоящая статья представляет собой пример документа, созданного с использованием языка разметки LaTeX. LaTeX является мощным инструментом для создания научных и технических документов, обеспечивая высокое качество форматирования текста и математических выражений. В этой статье мы рассмотрим основные элементы структуры документа и принципы его оформления.

LaTeX особенно популярен в академической среде благодаря своей способности корректно обрабатывать сложные математические формулы, библиографию и перекрестные ссылки. Его использование позволяет сосредоточиться на содержании документа, а не на его внешнем виде, поскольку стиль оформления определяется классом документа и пакетами.

\section{Основная часть}
В основной части статьи обычно размещается основной материал исследования или темы. В нашем примере это будет обзор основных возможностей LaTeX. Среди них - возможность создания таблиц, списков, вставки изображений и формул.

Например, простая формула может выглядеть так:
\begin{equation}
E = mc^2
\end{equation}

Также можно создавать списки:
\begin{itemize}
\item Первый пункт
\item Второй пункт
\item Третий пункт
\end{itemize}

\section{Заключение}
В заключение можно отметить, что LaTeX остается одним из наиболее эффективных инструментов для подготовки научных документов. Его использование позволяет создавать документы высокого качества с минимальными усилиями по форматированию. Благодаря своей стабильности и надежности, LaTeX продолжает оставаться выбором номер один для многих ученых и исследователей по всему миру.

В будущем планируется дальнейшее развитие и улучшение этого инструмента, чтобы удовлетворять растущие потребности научного сообщества в области подготовки и публикации научных материалов.

\end{document}